% ===================================================================
%  TAULUKKO N:o 9 — Vuori. — Kylpypaikka. — Metsä. — Metsästys.
%  Traduction 2026 d'apres texto/io, controlee sur texto/fr et
%  texto/sv. Consigne : voir texto/fi/10-taulukko-01.tex.
%
%  TROIS INVERSIONS EVITEES D'AVANCE, toutes de la meme forme :
%  « la sego (34) dil lignotranchisto (33) », « la hakilo (38) dil
%  dio-laboristo (39) », « la korno (53) dil chevalala chasisti (52) ».
%  Le possede porte un numero, le possesseur aussi, et le genitif
%  finnois les intervertirait. On passe donc par la relative :
%  « sahalle (34), joka on halonhakkaajalla (33) », et ainsi de suite.
% ===================================================================

%%K t09-apar-1 apar
\VUsaut{4.0mm}
\VUtitre{15.0pt}{60}{TAULUKKO N:o 9}
\VUsaut{3.0mm}

%%K t09-tit-1 sub
\VUcentre{12.0pt}{0}{\textit{Ensimmäinen kohtaus.}}
\VUsaut{3.0mm}

%%K t09-tit-2 sub
\VUpk{11.4pt}{40}{Vuori. --- Kylpypaikka.}
\VUsaut{3.0mm}

%%K t09-c1-01-1 p
1. --- Kahdeksan päivää olen ollut Pratsissa. En osaa ilmaista koko sitä
ihailua, jonka tunsin seisoessani Pyreneiden ihmeellisen \VUgras{vuorijonon}
\textsuperscript{(1)} edessä, jonka \VUgras{harja} \textsuperscript{(2)}
piirtyy siniselle taivaalle kuin jättiläismäinen köynnöskoriste.

%%K t09-c1-02-1 p
2. --- En kuvitellut, että Pyreneiden \VUgras{huiput} \textsuperscript{(3)}
saavuttavat niin suuren \VUgras{korkeuden}, ja seisoin hämmästyneenä komeiden
\VUgras{jäätiköiden} \textsuperscript{(5)} ja lumisten \VUgras{vuorenhuippujen}
\textsuperscript{(4)} edessä, joilla vierii niin hirvittäviä
\VUgras{lumivyöryjä}.

%%K t09-tit-3 sub
\VUsaut{3.0mm}
\VUpk{10.2pt}{40}{Vuoristomaisema.}
\VUsaut{3.0mm}

%%K t09-c1-03-1 p
3. --- Jo tuloni jälkeisenä päivänä menin vanhalle sammuneelle
\VUgras{tulivuorelle} \textsuperscript{(6)}, joka on Pratsin luona.
\VUgras{Kraatterin} kuivilla ja paljailla reunoilla näkee vielä hyvin jäljet
\VUgras{laavan} kulusta, joka virtasi useita vuosisatoja sitten pitkin
\VUgras{vuorenkylkeä} \textsuperscript{(7)}. Onnistuin, yhdellä
\VUgras{rinteistä}, löytämään muutamia palasia tuota laavaa.

%%K t09-c1-04-1 p
4. --- Prats on rajan luona, ja \VUgras{linnoitus} \textsuperscript{(8)}, joka
puolustaa \VUgras{vuorensolaa} \textsuperscript{(27)}, joka erottaa Ranskan
Espanjasta, muistuttaa aivan vanhoja linnoja Reinin rannalla, jotka kallionsa
huipulta näyttävät väijyvän \VUgras{saalista}.

%%K t09-c1-05-1 p
5. --- Valtava \VUgras{kotka} \textsuperscript{(9)}, jolla on pesänsä lähellä
\VUgras{linnaketta} \textsuperscript{(8)}, kiinnittää usein huomioni. Tuo
\VUgras{petolintu}, jonka \VUgras{nokka} on vahva, vie usein poikasilleen
karitsan tai kilin, jota se pitää \VUgras{kynsillään}. Niin suuri on sen
\VUgras{siipien} leveys, että se voi kauan liitää raskaine taakkoineen.

%%K t09-tit-4 sub
\VUsaut{3.0mm}
\VUpk{10.2pt}{40}{Vuorikävelijä.}
\VUsaut{3.0mm}

%%K t09-c1-06-1 p
6. --- Siellä miellyttävin kävelyretki on retki vesiputoukselle
\textsuperscript{(10)} « Cau'n » luona. \VUgras{Vesiputous} syöksyy pyörteinä
alas ja katoaa sitten kauniina \VUgras{purona} \VUgras{kuusimetsään}
\textsuperscript{(11)} kaupungista länteen. Nousimme tuon metsän läpi aina
\VUgras{ilmatieteelliselle havaintoasemalle} \textsuperscript{(12)} asti, ja
\VUgras{näköalapaikan} huipulta näimme koko seudun \VUgras{panoraaman}.
\VUgras{Maisema} on ihmeellinen, ja \VUgras{luonto} näyttää tuhlaavan maalle
kaikki ne aarteet, joita sillä on käytettävissään.

%%K t09-c1-07-1 p
7. --- Laskeutuaksemme käytimme \VUgras{vuorirataa} \textsuperscript{(13)} ja
pysähdyimme pienelle \VUgras{asemalle} vanhan \VUgras{feodaalilinnan}
\VUgras{raunioiden} \textsuperscript{(14)} viereen, jossa muinoin asuivat
Pratsin herrat.

%%K t09-c1-08-1 p
8. --- Poloinen linna! Mitä se mahtaakaan ajatella, kun se alapuolellaan näkee
keimailevan \VUgras{kylpypaikan} \textsuperscript{(15)}, joka nyt on
\VUgras{nykyaikainen parantola}?

%%K t09-c1-08-2 p
\VUgras{Ilmasto} on niin miellyttävä, \VUgras{kivennäisvesi} niin tehokas, että
\VUgras{hoito}, jota käytetään \VUgras{parantolassa} \textsuperscript{(17)}, on
todellinen nautinto.

%%K t09-tit-5 sub
\VUsaut{3.0mm}
\VUpk{10.2pt}{40}{Miellyttävä kylpykaupunki.}
\VUsaut{3.0mm}

%%K t09-c1-09-1 p
9. --- Muutoinkin on tehty kaikki oleskelun tekemiseksi miellyttäväksi. Suuri
\VUgras{maasilta} \textsuperscript{(16)} rakennettiin rautatien
mahdollistamiseksi; \VUgras{kylpylaitoksen} \VUgras{puisto}
\textsuperscript{(18)}, jossa annetaan \VUgras{konsertteja}, on istutettu
hurmaavalla tavalla. Useita sieviä « \VUgras{huviloita} »
\textsuperscript{(19)} rakennettiin kasinon \textsuperscript{(21)} ympärille, ja
hyvin hoidettu \VUgras{maantie} \textsuperscript{(22)} helpottaa yhteyksiä
suuresti.

%%K t09-c1-10-1 p
10. --- Pelkkä \VUgras{luiska} \textsuperscript{(23)} erottaa \VUgras{tien}
\textsuperscript{(20)} varsinaisesta \VUgras{laaksosta} \textsuperscript{(25)},
jonka pohjalla on sievä pieni \VUgras{järvi} \textsuperscript{(26)}, joka
ruokkii kirkasta \VUgras{puroa} \textsuperscript{(24)}.

%%K t09-c1-11-1 p
11. --- Laakson toisella puolella louhitaan \VUgras{rautakaivosta}
\textsuperscript{(28)}. Useaan kertaan menin katsomaan \VUgras{kaivosmiesten}
laskeutumista \VUgras{kuiluun}, ja menin jopa \VUgras{käytävään}, jossa he
viettävät päivänsä louhiessaan \VUgras{malmia}, joka sisältää
\VUgras{metallia}.

%%K t09-c1-12-1 p
12. --- Tärkeä \VUgras{rautatehdas} rakennettiin kaivoksen luo hyödyntämään heti
niitä rikkauksia, jotka siitä louhitaan. \VUgras{Masuuni}
\textsuperscript{(29)}, kuumennettuna sillä \VUgras{kivihiilellä}, jota täällä
on runsaasti, tekee mahdolliseksi saada nopeasti ja halvalla
\VUgras{takkirautaa}.

%%K t09-c1-13-1 p
13. --- Kaivoksesta ei kaukana kaivetaan \VUgras{louhosta}
\textsuperscript{(30)}. Ei \VUgras{graniittia}, eikä \VUgras{porfyyriä}, eikä
edes \VUgras{hiekkakiveä} hakkaa vanha \VUgras{kivenhakkaaja}
\VUgras{hakullaan}, vaan yksinkertaista \VUgras{hakattua kiveä} talonrakennukseen.

%%K t09-c1-14-1 p
14. --- Yksi tuon maan kauneimmista koristeista ovat \VUgras{kuuset}
\textsuperscript{(31)}. Kaikkialla \VUgras{rotkon} \textsuperscript{(32)}
pohjalla, \VUgras{kosken} \textsuperscript{(33)} rannalla,
\VUgras{kuilun} \textsuperscript{(34)} tai jyrkänteen partaalla, niiden ohuet
neulaset vivahtavat vihreään.

%%K t09-tit-6 sub
\VUsaut{3.0mm}
\VUpk{10.2pt}{40}{Jalan kulkeminen.}
\VUsaut{3.0mm}

%%K t09-c1-15-1 p
15. --- Käytän lomani \VUgras{pitkiin} kävelyretkiin. \VUgras{Tiet}
\textsuperscript{(35)}, jotka kiemurtelevat vuorelle, tekevät mahdolliseksi,
välittämättä etäisyydestä, taittaa useita \VUgras{kilometrejä} ja usein useita
\VUgras{peninkulmia}.

%%K t09-c1-15-2 p
\VUgras{Rajakivet} \textsuperscript{(36)} ja \VUgras{tienviitat}
\textsuperscript{(37)} merkitsevät tiet, joilla usein kohtaa nuoren
\VUgras{vuoristolaisen} \textsuperscript{(38)}, jolla on \VUgras{reppu}
\textsuperscript{(40)} kyljellään ja joka kantaa \VUgras{murmelia}
\textsuperscript{(39)}, tai jonkin \VUgras{mustalaisen} \textsuperscript{(41)},
jonka \VUgras{turkislakki} \textsuperscript{(42)} omalaatuisesti erottuu vanhasta
repaleisesta \VUgras{hupusta} \textsuperscript{(43)}. Hän taluttaa
\VUgras{ketjussa karhua} \textsuperscript{(44)}, joka on koulutettu.

%%K t09-tit-7 sub
\VUpk{10.2pt}{40}{Vuorikiipeily.}
\VUsaut{3.0mm}

%%K t09-c1-16-1 p
16. --- Miltei joka päivä menen \VUgras{oppaan} \textsuperscript{(48)} kanssa
tekemään \VUgras{nousun}, ja olen hyvin ylpeä, kun \VUgras{reppu}
\textsuperscript{(47)} selässäni ja « \VUgras{alppisauva} » kädessäni valloitan
kuin oikea \VUgras{turisti} \textsuperscript{(46)} \VUgras{kalliot}
\textsuperscript{(45)}.

%%K t09-c1-17-1 p
17. --- Vuoren huipulta tasangon asukkaat eivät näytä muurahaisia
suuremmilta; \VUgras{lammaslauma} \textsuperscript{(50)}, joka laiduntaa
\VUgras{laitumella} \textsuperscript{(49)}, näyttää lastenleikkikalulta.
\VUgras{Mänty} \textsuperscript{(51)} tai myös pieni \VUgras{puutalo}
\textsuperscript{(52)} näyttää tuskin täplältä vehreydessä.

%%K t09-c1-18-1 p
18. --- Noiden retkien aikana kohtaamme usein \VUgras{polulla}
\textsuperscript{(53)} \VUgras{gemssinmetsästäjän} \textsuperscript{(54)}, joka
kantaa olallaan eläintä, jonka on juuri surmannut.

%%K t09-c1-19-1 p
19. --- Panin kasvioni väliin hyvin merkillisen \VUgras{saniaisen}
\textsuperscript{(55)} lehden ja kauniin ruusunpunaisen \VUgras{kanervan}
\textsuperscript{(57)} oksan, jotka poimin pienen \VUgras{luolan}
\textsuperscript{(56)} suulta viimeisellä kävelyretkelläni.

%%K t09-tit-8 sub
\VUsaut{3.0mm}
\VUcentre{12.0pt}{0}{\textit{Toinen kohtaus.}}
\VUsaut{3.0mm}

%%K t09-tit-9 sub
\VUpk{11.4pt}{40}{Metsä. --- Metsästys.}
\VUsaut{3.0mm}

%%K t09-c2-01-1 p
1. --- Eilen vietin ihanan päivän metsässä. Se ahkeruus, joka vallitsee
pikkueläinten ja \VUgras{hyönteisten} \textsuperscript{(5)} keskuudessa, joiden
asuttama se on, kiinnosti minua kovasti.

%%K t09-c2-01-2 p
\VUgras{Hämähäkki} \textsuperscript{(1)}, joka kutoo \VUgras{verkkoaan}
\textsuperscript{(2)} kahden oksan väliin pyydystääkseen \VUgras{kärpäsen}
\textsuperscript{(3)}, joka lentää ohi, \VUgras{etana} \textsuperscript{(4)},
joka vaivalloisesti kantaa taloaan, sarvijaakko, joka etsii saalistaan, ahkera
muurahainen, joka hyvin innokkaasti puuhaa \VUgras{muurahaispesällä}
\textsuperscript{(6)}, kaikki nuo pienet menivät, tulivat, työskentelivät
tavattomalla ahkeruudella.

%%K t09-c2-02-1 p
2. --- \VUgras{Käärme} \textsuperscript{(7)}, jota ensi silmäyksellä pidin
\VUgras{kyynä}, mutta joka ei ollut mitään muuta kuin tavallinen
\VUgras{rantakäärme}, oli piiloutunut puunrungon alle, ja pisti silloin tällöin
esiin ohuen pienen päänsä, joka aina näytti valmiilta puremaan. Edessäni ryömi
paksu \VUgras{etana} \textsuperscript{(8)} raskaasti eteenpäin nieltyään yhden
niistä maukkaista \VUgras{ahomansikoista} \textsuperscript{(11)}, joita siellä
kasvaa runsaasti.

%%K t09-c2-03-1 p
3. --- Maa oli matoitettu \VUgras{metsäkukilla}; \VUgras{kevätesikot}
\textsuperscript{(9)}, \VUgras{talvio} \textsuperscript{(10)}, ja monet muut
kasvit, joita en tuntenut, kukkivat \VUgras{ruohotuppaiden}
\textsuperscript{(12)} välissä.

%%K t09-c2-04-1 p
4. --- Sillä seudulla \VUgras{tryffeli} on miltei yhtä tavallinen kuin
\VUgras{sieni} \textsuperscript{(14)}, ja \VUgras{porsas}
\textsuperscript{(13)}, joka kuului eräälle metsänvartijalle, etsi ahnaasti,
onnistuisiko se löytämään niitä muutamia.

%%K t09-tit-10 sub
\VUsaut{3.0mm}
\VUpk{10.2pt}{40}{Salametsästys.}
\VUsaut{3.0mm}

%%K t09-c2-05-1 p
5. --- Olin läsnä erään kohtauksen \VUgras{lopussa}, joka huvitti minua kovasti.
\VUgras{Mäyräkoiransa} \textsuperscript{(15)} hajuaistin avulla
\VUgras{metsänvartija} \textsuperscript{(16)} oli juuri yllättänyt
\VUgras{salametsästäjän} \textsuperscript{(22)} sillä hetkellä, kun tämä
valmistautui kantamaan pois \VUgras{riistan} \textsuperscript{(17)}, jonka oli
surmannut. Koska hän mieluummin hylkäsi saaliinsa kuin antoi ottaa itsensä
kiinni, salametsästäjä oli paennut, ja \VUgras{jänis} \textsuperscript{(18)},
\VUgras{naarashirvi} \textsuperscript{(19)}, \VUgras{metsäkauris}
\textsuperscript{(20)} ja \VUgras{villisika} \textsuperscript{(21)} makasivat
ruohikolla metsänvartijan iloksi.

%%K t09-c2-06-1 p
6. --- Niiden puiden moninaisuus, joita metsä sisältää, on hämmästyttävä.
\VUgras{Pyökit} \textsuperscript{(23)} ja \VUgras{koivut}
\textsuperscript{(26)}, \VUgras{männyt}, jotka antavat \VUgras{pihkan}
\textsuperscript{(27)}, \VUgras{tammet} \textsuperscript{(29)} ja myös monet
muut sekoittavat oksistonsa.

%%K t09-c2-06-2 p
\VUgras{Muratti} \textsuperscript{(24)}, joka takertuu korkeiden puiden runkoon,
värjää \VUgras{metsän} \textsuperscript{(25)} kiiltävän vihreäksi, joka
miellyttävästi yhtyy \VUgras{sammalen} \textsuperscript{(31)} vaaleampaan ja
kalpeanvihreään värisävyyn, jossa \VUgras{harmaapäätikka}
\textsuperscript{(30)} etsii hyönteisiä.

%%K t09-tit-11 sub
\VUsaut{3.0mm}
\VUpk{10.2pt}{40}{Metsämaisema.}
\VUsaut{3.0mm}

%%K t09-c2-07-1 p
7. --- Vanhat tammet lumosivat minut. Niiden paksut kiertyneet \VUgras{juuret}
\textsuperscript{(32)}, puoliksi maasta nousseina, antavat niille komean voiman
ja väkevyyden ilmeen, ja kysyn itseltäni, kuinka \VUgras{omistaja}
\textsuperscript{(35)} voi saada itsensä kaatamaan ne ja luovuttamaan ne
\VUgras{sahalle} \textsuperscript{(34)}, joka on \VUgras{halonhakkaajalla}
\textsuperscript{(33)}. Poloiset \VUgras{rungot} \textsuperscript{(36)}, joilta
on riistetty \VUgras{kuori} \textsuperscript{(37)} tai jotka on halkaistu
\VUgras{kirveellä} \textsuperscript{(38)}, joka on \VUgras{päivätyöläisen}
\textsuperscript{(39)} kädessä, surettavat minua.

%%K t09-c2-08-1 p
8. --- Aivan lähellä vanha \VUgras{sysimies} \textsuperscript{(41)} oli
\VUgras{lapiollaan tehnyt kasan} \textsuperscript{(42)} hiiltä, ja vanha nainen
kantoi \VUgras{kimppua} \textsuperscript{(40)} \VUgras{kuivia puita}
kaadettujen puiden oksista.

%%K t09-tit-12 sub
\VUsaut{3.0mm}
\VUpk{10.2pt}{40}{Metsästys.}
\VUsaut{3.0mm}

%%K t09-c2-09-1 p
9. --- Halusin mennä lepäämään muutamaksi hetkeksi \VUgras{metsänvartijan taloon}
\textsuperscript{(46)}, mutta, kun tulin \VUgras{pienelle järvelle}
\textsuperscript{(43)}, jolla kaunis \VUgras{joutsen} \textsuperscript{(45)} ui,
huomasin, että äskettäin sattunut \VUgras{tulva} \textsuperscript{(44)} oli
muuttanut tien todelliseksi \VUgras{suoksi}. Minun täytyi poiketa syrjään, ja,
kun kuljin \VUgras{setrin} \textsuperscript{(49)}, \VUgras{poppelien}
\textsuperscript{(48)} ja \VUgras{jalavan} \textsuperscript{(47)} ohi, jotka ovat
metsän reunassa, näin tulevan esiin \VUgras{pensaikosta}
\textsuperscript{(54)} komean \VUgras{hirven} \textsuperscript{(50)}, jota
ajoi takaa itsepintainen \VUgras{koiralauma} \textsuperscript{(51)}, jota myös
yllytti \VUgras{torvi} \textsuperscript{(53)}, jota \VUgras{ratsastavat
metsästäjät} \textsuperscript{(52)} puhalsivat. Poloinen eläin oli aivan uuvuksissa.
Se tuli ja kaatui kuolemaisillaan muutaman askeleen päähän minusta.

%%K t09-c2-10-1 p
10. --- Metsänvartijan poika, jonka \VUgras{ajometsästyksen} melu oli
houkutellut sinne, tuli ulos talosta, ja menimme yhdessä takaisin metsään. Hän
oli nähnyt \VUgras{harakan} \textsuperscript{(56)} vanhan tammen oksalla. Hän
arveli, että sen pesä on luultavasti kätketty \VUgras{lehvistöön}
\textsuperscript{(58)}, ja hän tahtoi ottaa pesän.

%%K t09-c2-10-2 p
Ketterästi kuin \VUgras{orava} \textsuperscript{(61)} hän kiipesi
\VUgras{oksalta} \textsuperscript{(55)} oksalle, mutta hän ei löytänyt mitään ja
onnistui vain säikäyttämään vanhan \VUgras{varislinnun} \textsuperscript{(60)},
joka istui \VUgras{paksulla oksalla} \textsuperscript{(57)}.
\VUsaut{4.0mm}
\VUfilet{22mm}
